\documentclass{article}
\input{../../../../../newpreamble.tex}
\title{Homework 5 Math 210A}
\author{Lance Remigio}
\begin{document}
\maketitle

\begin{problem}
    Let \( \varphi: {G}_{1} \to {G}_{2} \) be a group homomorphism.
    \begin{enumerate}
        \item[(a)] Show that if \( {N}_{2} \trianglelefteq {G}_{2}  \), then \( \varphi^{-1}({N}_{2}) \trianglelefteq {G}_{1}  \). Deduce that \( \text{ker}(\varphi) \trianglelefteq {G}_{1} \).
            \begin{proof}
            Suppose \( {N}_{2} \trianglelefteq {G}_{2} \). Our goal is to show that \( \varphi^{-1}({N}_{2}) \trianglelefteq {G}_{1} \); that is, for all \( g \in {G}_{1} \),  
            \[  g \varphi^{-1}({N}_{2}) g^{-1} \subseteq \varphi^{-1}({N}_{2}). \tag{*} \]
            To this end, let \( g \in {G}_{1} \) be arbitrary. To show (*), let \( x \in g \varphi^{-1}({N}_{2}) g^{-1} \) be arbitrary. Then 
            \[  x = g \tilde{x} g^{-1} \ \ \text{for some} \  \tilde{x} \in \varphi^{-1}({N}_{2}). \]
            Since \( \tilde{x} \in \varphi^{-1}({N}_{2}) \), it follows that \( \varphi (\tilde{x}) \in {N}_{2}  \). But note that \( {N}_{2} \trianglelefteq {G}_{2} \) and so for all \( g' \in {G}_{2} \), we have  
            \[  g' {N}_{2} (g')^{-1}  = {N}_{2}. \tag{\( \dagger \)} \]
            Since \( \varphi  \) is a homomorphism, we see that 
            \begin{align*}
                \varphi(x) &= \varphi(g \tilde{x} g^{-1}) 
                           = \varphi(g) \varphi(\tilde{x}) (\varphi(g))^{-1}    
            \end{align*}
            Since \( \varphi(g) \in {N}_{2} \), we have \( \varphi(x) \in {N}_{2} \) from (\( \dagger \)). Thus, \( x \in \varphi^{-1}({N}_{2}) \) and so we conclude that   
            \[  g \varphi^{-1}({N}_{2})g^{-1} \subseteq  \varphi^{-1}({N}_{2})\]
            which is our desired result. We see that \( \text{ker}(\varphi) \trianglelefteq {G}_{1} \) because
            \begin{align*}
                \text{ker}(\varphi) &= \{  g \in {G}_{1} : \varphi(g) = {1}_{G} \varphi^{-1}({N}_{2}) \}  \\
                                    &= \{  g \in {G}_{1} : g {N}_{2} =  \varphi^{-1} ({N}_{2})  \}  \\
                                    &= \{  g \in {G}_{1} : g \in \varphi^{-1}({N}_{2}) \} \\ 
                                    &= \varphi^{-1}({N}_{2}) \\
                                    &\trianglelefteq {G}_{1}.
            \end{align*}
            \end{proof}
        \item[(b)] Show that if \( {N}_{1} \trianglelefteq {G}_{1} \), then \( \varphi({N}_{1}) \trianglelefteq \varphi({G}_{1}) \). Deduce that any surjective homomorphism will map normal subgroups in the domain to normal subgroups in the codomain.
            \begin{proof}
            Suppose that \( {N}_{1} \trianglelefteq {G}_{1} \). To show that \( \varphi({N}_{1}) \trianglelefteq \varphi({G}_{1}) \), we need to show that for all \( g' \in \varphi({G}_{1}) \), 
            \[ g' \varphi({N}_{1}) (g')^{-1} \subseteq  \varphi({N}_{1}). \tag{*}  \]
            Let \( g' \in \varphi({G}_{1}) \). To show (*), let \( y \in g' \varphi({N}_{1}) (g')^{-1} \). Then for some \( \tilde{y} \in \varphi({N}_{1})\), we have 
            \[  y = g' \tilde{y} (g')^{-1}. \]
            Since \( \tilde{y} \in \varphi({N}_{1})  \), then for some \( \tilde{x} \in {N}_{1} \), we have \(  \tilde{y} = \varphi(\tilde{x}) \). Furthermore, \( g' \in \varphi({G}_{1}) \) implies that there exists some \( \tilde{g} \in  {G}_{1} \) such that \( g' = \varphi(\tilde{g}) \). Since \( \varphi  \) is a homomorphism, it follows that 
            \begin{align*}
                y = g' \tilde{y} (g')^{-1} &= \varphi(\tilde{g}) \varphi(\tilde{x}) \varphi(\tilde{g})^{-1} \\
                                           &= \varphi(\tilde{g}) \varphi(\tilde{x}) \varphi(\tilde{g}^{-1}) \\
                                           &= \varphi(\tilde{g} \tilde{x} \tilde{g}^{-1}).
            \end{align*}
            Notice that \( \tilde{g} \tilde{x} \tilde{g}^{-1} \in \tilde{g} {N}_{1} \tilde{g}^{-1}  \). But since \( {N}_{1} \trianglelefteq {G}_{1} \), we have for all \( g \in G  \), \( g {N}_{1} g^{-1} = {N}_{1} \). In particular, \( \tilde{g} {N}_{1} \tilde{g}^{-1} = {N}_{1} \). Hence,  
            \( \hat{x} = \tilde{g} \tilde{x} \tilde{g}^{-1} \in {N}_{1}  \). Hence, we can see that \( y \in \varphi({N}_{1}) \) and so we conclude that 
            \[ g' \varphi({N}_{1}) (g')^{-1} \subseteq  \varphi({N}_{1}).   \]
            Hence, we conclude that \( \varphi({N}_{1}) \trianglelefteq \varphi({G}_{1}) \).
            \end{proof}
    \end{enumerate}
\end{problem}

\begin{problem} Let \( G  \) be a group with \( N \trianglelefteq G  \).
    \begin{enumerate}
        \item[(a)] Suppose \( x \in G  \) with \( |  x  |  = n < \infty   \). Show that \( n \) divides the order of \( x N  \) in \( G / N  \).
            \begin{proof}
            Our goal is to show that \( n \mid | x N  |  \). Denote \( | x N  | = k   \). By the division algorithm, we can find a unique \( q \in \Z  \) and \( 0 \leq r < n  \) such that  
            \[  k = n q + r. \tag{*} \]
            Since \( N \) is a normal subgroup, we know that coset multiplication in \( G / N  \) is well-defined. Hence, we see that 
            \[ x^{k} = x^{n q  +  r} = (x^{n})^{q} \cdot x^{r} = x^{r}.\]
            Note that \( r = 0  \) because otherwise if \( r \neq 0   \) such that \( x^{r} = 1  \), then \( r  \) would contradict the minimality of \( n  \). Hence, from (*) we can see that \( | x  |  = n  \) must divide \( k  \). Thus, \( n  \) divides the order of \( x N  \) in \( G / N  \).   
            \end{proof}
        \item[(b)] Further show that the order \(  xN  \) in \( G / N  \) is the smallest positive integer \( k  \) such that \( x^{k } \in N  \).
            \begin{proof}
                Suppose \( N \neq \{ 1_G  \}  \). Then there exists some \( a \neq 0  \) such that \( x^{a} \in N  \). If \( a < 0  \), then since \( N  \) is a group \( x^{-a} = (x^{a})^{-1} \in N \). This tells us that the set \( \mathcal{P} = \{  b \in \Z^{+} : x^{b} \in N \}  \) is nonempty. Since \( \mathcal{P} \) is a subset of \( \Z  \), the well-ordering property implies that there exists a minimal element \( d  \) such that \( x^{d} \in  N \). In \( G / N  \), this corresponds to 
    \[  x^{d} N = {1}_{G} N.  \]
    But since the order of \( xN  \) in \( G / N  \) is unique, this tells us that \( k = d  \). Hence, \( k  \) is the smallest positive integer such that \( x^{k} \in N  \).
            \end{proof}
        \item[(c)] Show that if \( G  =  \langle x   \rangle  \), then \( G / N  \) is cyclic.
            \begin{proof}
            Since \( G = \langle x  \rangle  \) and \( | x  |  = n  \). Then \( | \langle x  \rangle |  = | G | = n   \).   Our goal is to show that \( \langle x N  \rangle = G / N  \). From parts (a) and (b), we can see that \( | x N  |  = k \) and \( k  \) is the smallest such positive integer such that \( x^{k} \in N  \); that is, \( (xN)^{k} \in G / N  \). Hence, we see that \( \langle x N  \rangle \leq G / N  \). All that remains to show is that \( G / N \leq \langle x N  \rangle\). Let \( t  \) be any positive power of \( xN \). We will show that \( (xN)^{t} \in \langle x N  \rangle \). Using the division algorithm, we can find a unique \( q \in \Z  \) and \( 0 \leq r < k  \) such that 
                \[  t = k q + r.  \]
                Since coset multiplication is well-defined in \( G /N  \), we have that 
                \begin{align*}
                    (xN)^{t} = x^{t} N  = (x^{kq + r}) N  
                                        &= \Big(  (x^{k})^{q} \cdot x^{r} \Big) N \\ 
                                        &= x^{r} N   \tag{\( x^{k} = 1  \) from part (a)} \\
                                        &= (xN)^{r}. \tag{\( 0 \leq r < k \)}
                \end{align*}
                Hence, \( (xN)^{t} \in \langle x N   \rangle  \) and so we have \( \langle x N   \rangle = G / N  \).
            \end{proof}
    \end{enumerate}
\end{problem}

\begin{problem}
    Let \( G  \) with center \( Z(G) \). Recall that \( Z(G) \trianglelefteq G  \).
    \begin{enumerate}
        \item[(a)] Prove that if \( G / Z(G) \) is cyclic, then \( G  \) is abelian.
            \begin{proof}
          Our goal is to show that \( G  \) is abelian. Suppose \( G / Z(G) \) is cyclic. We need to show that 
          \[  xy = yx \ \ \forall x,y \in G. \]
          To this end, let \( x,y \in G  \) be arbitrary. Since \( G / Z(G) \) is cyclic, we must have that \( G / Z(G) \) is abelian. Since \( xy \in G  \) and \( yx \in G  \), we have that (because \( Z(G) \trianglelefteq G  \) and so the coset multiplication is well-defined) 
          \[  (xy) Z(G) = (x Z(G) )(y Z(G)) = (y Z(G))(x Z(G)) = (yx) Z(G).  \]
          This tells us that \( (xy)^{-1}(yx) \in Z(G) \) and so we have 
          \[ (xy)^{-1}(yx)g = g (xy)^{-1}(yx) \ \ \forall g \in G.  \]
          In particular, if \( g = x^{-1} \), we have 
          \begin{align*}
              &(xy)^{-1}(yx) x^{-1} = x^{-1}(xy)^{-1} (yx) \\
              &\implies y^{-1}x^{-1} y = x^{-1}y^{-1} x^{-1}yx \\    
              &\implies x(y^{-1} x^{-1} y) = (y^{-1}x^{-1}y)x.
          \end{align*}
          Since \( x \in G  \) is arbitrary, we see that \( y^{-1}x^{-1} y \in Z(G) \). If \( g = y  \), then
          \begin{align*}
              &y(y^{-1}x^{-1}y) = (y^{-1}x^{-1}y)y \\
              &\implies x^{-1}y = (y^{-1}x^{-1}y) y \\
              &\implies x^{-1} = y^{-1}x^{-1}y \\
              &\implies xy = yx.
          \end{align*}
          Hence, \( G  \) must be abelian.
            \end{proof}
        \item[(b)] Show that if \( | G  |  = pq  \) where \( p  \) and \( q  \) are primes (not necessarily distinct) then either \( G  \) is abelian or \( Z(G) = \{  {1}_{G} \}  \).
            \begin{proof}
            Suppose \( | G  |  = pq  \) where \( p  \) and \( q  \) are primes. Suppose \( Z(G) \neq \{ {1}_{G} \}  \). Our goal is to show that \( G  \) is abelian. Since \( G  \) is a finite group and \( Z(G) \leq G  \), it follows from Lagrange's Theorem that  
            \[  | Z(G) | \mid | G |. \]
            This tells us that \( | Z(G) | \in \{ 1, p, q , pq  \}   \). If \( | Z(G)  |  = 1  \), then we would get a contradiction because we assumed that \( Z(G) \neq 1  \). If \( | Z(G) | = p  \), then 
            \[  | G / N  |  = \frac{ | G |  }{ | Z(G) |  }  = \frac{ pq  }{ p  } = q. \]
            Since \( q  \) is prime, \( G / N \) must be a cyclic group of order prime \( q  \). From part (a), \( G  \) must be abelian. Similarly, if \( | Z(G) |  = q  \), then we would get \( | G / N  |  = p   \), and so \( G / N  \) is a cyclic group of order prime \(  p \) and hence \( G  \) is an abelian group. If \( | Z(G) |  = pq = | G |   \) and \( Z(G) \leq G  \), then \( Z(G) = G   \) and we also get that \( G  \) is abelian.

            If \( p  \) and \( q  \) are not distinct, then \( | Z(G) | \in \{ 1, p, p^{2} \}  \). A similar argument will show that \( G  \) will be an abelian group given that \( | G / N |  \) is prime. For the case that \( | Z(G) | = p^{2}  \), then \( Z(G) = G  \) and so \( G  \) is abelian. 
            \end{proof}
    \end{enumerate}
\end{problem}

\begin{problem}[Section 3.2 Problem 5]
  Let \( H  \) be a subgroup of \( G  \) and fix some element \( g \in G  \).  
  \begin{enumerate}
      \item[(a)] Prove that \( g H g^{-1}  \) is a subgroup of \( G  \) of the same order as \( H  \).
          \begin{proof}
          \end{proof}
        \item[(b)] Deduce that if \( n \in \Z^{+} \) and \( H  \) is the unique subgroup of order \( n  \) then \( H \trianglelefteq G  \).
            \begin{proof}
            \end{proof}
  \end{enumerate}
\end{problem}

\begin{problem}[Section 3.2 Problem 8]
  Prove that if \( H  \) and \( K  \) are finite subgroups of \( G  \) whose orders are relatively prime then \( H \cap  K = 1  \).  
  \begin{proof}
  
  \end{proof}
\end{problem}


\end{document}

